\section{\label{sec:ope}The operator product expansion}

Wilson's approach to the OPE for a nonlocal operator is widely known---see, for 
example, Ref.~\cite{Collins:1984rdg}---and here we review some notation 
necessary for our discussion.

We write the OPE for a nonlocal operator, ${\cal Q}(x)$, as
\begin{equation}
{\cal Q}(x)  \stackrel{x\rightarrow 0}{\sim} \; \sum_k 
c_k(x,\mu) [{\cal O}^{(k)}(0)]_{\mathrm{R}}.
\end{equation}
The Wilson coefficients $c_k(x,\mu)$ are complex functions of the spacetime 
separation, $x$, and renormalization scale, $\mu$, that capture the 
short-distance physics associated with the renormalized local operator $[{\cal 
O}^{(k)}(0)]_{\mathrm{R}}$. We represent renormalized operators by 
$[\ldots]_{\mathrm{R}}$ and suppress their dependence on the renormalization 
scale, $\mu$, for notational simplicity. The 
free-field mass dimension of the local operator determines the leading 
spacetime dependence of the corresponding Wilson coefficient.

As an example, let us consider the time-ordered two-point function of two 
scalar fields with spacetime separation $x$: ${\cal T} \{\phi(x)\phi(0)\}$. In 
free scalar field theory, 
the OPE is a Laurent expansion. Interactions generate subleading dependence 
on the spacetime separation in the Wilson coefficients, which become 
functions of the spacetime separation, the (renormalized) mass 
$m_{\mathrm{R}}$, and the 
renormalization scale, $\mu$:
\begin{align}\label{eq:opeint}
{\cal T} \{\phi(x)\phi(0)\} = {}& \frac{c_{\mathbb{I}}(\mu 
x,m_{\mathrm{R}}x)}{4\pi^2x^2}\mathbb{I}+ 
c_{\phi^2}(\mu x,m_{\mathrm{R}}x)[\phi^2(0) ]_{\mathrm{R}} \nonumber\\
{} & \qquad 
+{\cal O}(x).
\end{align}
Here we have 
factored 
out 
the leading spacetime dependence from the Wilson 
coefficients. The ${\cal O}(x)$ indicates terms of order $x$, up to 
logarithmic corrections generated by interactions.

We propose replacing the set of local operators in the OPE by
their locally smeared counterparts
\begin{equation}
{\cal Q}(x)  \stackrel{x\rightarrow 0}{\sim} \; \sum_k 
d_k(\tau,x,\mu) {\cal S}^{(k)}(\tau,0).
\end{equation}
The Wilson coefficients, $d_k(\tau,x,\mu)$, and the smeared 
operators, ${\cal S}^{(k)}(\tau,0)$, are now functions of an extra 
scale, the 
flow time, $\tau$. Nevertheless, the leading spacetime 
dependence of the Wilson coefficients is dictated by the canonical mass 
dimension of the corresponding smeared operator
and is therefore just that of the standard OPE. 

Just as the OPE is only valid 
for small spacetime separations, the sOPE requires small flow times. We will 
see that we require that $\tau \propto x^2$, where $x$ is assumed to be small, 
to ensure that the Wilson coefficients are independent of the external states.

For example, returning to the time-ordered two-point 
function, the sOPE is
\begin{align}\label{eq:sope2pt}
\!\!{\cal T}\{\phi(x)\phi(0)\} ={} &
\frac{d_{\mathbb{I}}(\mu x,\mu^2\tau,m_{\mathrm{R}}x)}{4\pi^2x^2}\mathbb{I}
\nonumber\\
{} & + 
d_{\rho^2}(\mu x,\mu^2\tau,m_{\mathrm{R}}x)\rho^2(\tau,0)  + 
{\cal O}(x,\tau),
\end{align}
where we denote smeared fields at flow time $\tau$ by $\rho(\tau,x)$.

In the following sections, we will observe four features of the 
smeared Wilson coefficients in the sOPE.
\begin{enumerate}
\item The logarithmic spacetime dependence of the original OPE is preserved in 
the sOPE.
\item The flow time serves as an ultraviolet regulator for the smeared Wilson 
coefficients, to leading-order in perturbation theory. 
\item Beyond leading-order the flow time cannot regularize the Wilson 
coefficients, because the flow evolution is classical. We can, however, absorb 
renormalization scale dependence into the renormalization parameters 
of the original theory.
\item For any OPE to be meaningful, the Wilson coefficients must be independent 
of the external states. We can ensure this for the sOPE by choosing 
$s_{\mathrm{rms}} < x$; \emph{i.e.}~the mean smearing radius must be smaller 
than the spacetime extent of the nonlocal operator. This ensures that the sOPE
remains an expansion in approximately local operators. In other words, if the 
gradient flow probes length scales on the order of the nonlocal operator, then 
the sOPE becomes a poor expansion for the original operator.
\end{enumerate}

Although we refer to this expansion as the \emph{smeared} OPE, we really have 
in mind that the smearing is implemented via the gradient flow. The 
gradient flow acts as a smoothing operation that
drives the degrees of freedom of the theory to the stationary points of the 
action. In
QCD, the gradient flow corresponds to a continuous stout-smearing procedure, an 
analytic method for constructing lattice gauge fields with damped ultraviolet
fluctuations \cite{Morningstar:2003gk}. The direct comparison of the gradient 
flow to other smearing schemes was first undertaken in 
Ref.~\cite{Bonati:2014tqa}.

Many studies of the gradient flow have incorporated a small flow-time expansion 
of fields at nonvanishing flow time in terms of local 
fields at zero 
flow time \cite{Shindler:2013bia,Suzuki:2013gza,DelDebbio:2013zaa, 
Makino:2014taa,Asakawa:2013laa,Luscher:2014kea}. We can view such an expansion 
as an OPE in the flow time and thereby relate 
renormalized quantities calculated at nonzero flow time to the corresponding 
quantities in the original theory at vanishing flow time, which would otherwise 
be difficult to compute.

In this work, we take a different approach. We do 
not expand the flowed fields in terms of original fields at nonzero flow time, 
but rather take 
as the fundamental objects of study the (matrix elements of)
fields at positive flow time. Both approaches reflect the physically-motivated 
expectation that smearing 
scales much smaller than the physical scales of the system should not distort 
the physics in question. The small flow-time expansion quantifies any 
deviations 
from this expectation and, 
furthermore, decouples analytic calculations of Wilson coefficients in the 
continuum from lattice calculations with smeared degrees of freedom. The sOPE 
incorporates the smearing scale as an inherent scale of the system, which 
requires new Wilson coefficients to be determined. Thus, both the small 
flow-time 
expansion and the sOPE are shaped by the role of the smearing scale as 
ultraviolet regulator and are related, but distinct, conceptual approaches to 
the same physics: partially restoring rotational symmetry on the lattice.

Smearing, in 
general, is a tool that partially restores rotational symmetry 
\cite{Davoudi:2012ya} and thereby improves the continuum limit of lattice 
calculations. Smearing via the gradient flow has a number of advantages. In 
particular, 
matrix elements determined nonperturbatively on the lattice using 
smeared degrees of freedom require no further renormalization 
\cite{Luscher:2010iy}, up to fermionic 
wave function renormalization \cite{Luscher:2013cpa}, and remain 
finite, provided the smearing scale is kept fixed in physical units as the 
continuum limit is taken.

We now turn to a more complete study of the sOPE applied to $\phi^4$ scalar 
field theory, a particularly straightforward theory in which to examine 
the sOPE. We can solve the flow-time 
equations exactly, because the flow-time 
evolution is linear in the scalar field.
We do not need to consider the complications associated with 
gauge fixing \cite{Luscher:2010iy,Luscher:2011bx}---nor the extra 
renormalization of fermions \cite{Luscher:2013cpa}. Furthermore, for scalar 
fields the sOPE can be understood to be simply a resummation of the original 
OPE. Although it is not necessary for our work, it is also interesting to note 
that 
the OPE is known to converge for Euclidean $\phi^4$ theory in four dimensions, 
at a fixed, but arbitrary order in the perturbative 
expansion \cite{Hollands:2011gf}. This result holds at arbitrary spacetime 
separations, provided the external states are of compact support.

Looking toward future calculations in QCD, we anticipate that 
the technical issues associated with a flow equation that incorporates gauge 
field interactions will slightly complicate the perturbative calculations by 
increasing the number of diagrams at a given order in perturbation theory, but 
will not alter our conclusions. Within the gauge sector, there are no loops of 
flowed fields, because renormalized correlation functions remain finite at 
nonzero flow time \cite{Luscher:2011bx}. Therefore, at leading-order in 
perturbation theory, the flow time regulates ultraviolet divergences; beyond 
leading-order an appropriate renormalization procedure must be incorporated. 
With fermionic fields there is an extra fermion renormalization parameter, 
calculated in Ref.~\cite{Luscher:2013cpa}, but this can be removed by 
considering 
renormalization group invariant quantities \cite{Luscher:2013cpa}. We also note 
that all perturbative calculations for the sOPE can be carried out in the 
continuum and do not require lattice perturbation theory, which is generally 
more involved \cite{Monahan:2013dla,Capitani:2002mp}.
